\section{Introduction}
\label{sec:introduction}
In this paper, we study the structured nonsmooth nonconvex composite problem
\begin{align}
    \min_{\bx\in\RR^n}\quad & F(\bx) := f(\bx) + h(\calA(\bx)),
    \tag{\textbf{P}}
    \label{prob:main}
\end{align}
where $f\colon\RR^n\to\RR$ is a possibly nonconvex function with a Lipschitz
continuous gradient, $\calA\colon\RR^n\to\RR^m$ is a smooth and possibly nonlinear
mapping, and $h\colon\RR^m\to\RR\cup\{+\infty\}$ is a proper closed convex function
whose proximal mapping is easy to evaluate; the complete assumptions are stated in
Section~\ref{sec:assumptions}. The two components play complementary roles: $f$ is
usually a smooth loss or data-fidelity term, whereas $h$ is a nonsmooth convex function
that imposes robustness, sparsity, or other structure on the transformed variable
$\calA(\bx)$
\cite{charisopoulos2021low,markovsky2008structured,parekh2015convex}. Several models
of practical interest take this form.  
\subsection{Motivating Examples}
\textbf{Robust phase retrieval}~\cite{duchi2019solving}.
Given measurement vectors $\mathbf a_i\in\RR^n$ and observations
$b_i\approx\langle\mathbf a_i,\bx^\star\rangle^2$, a fraction of which may be grossly
corrupted, one solves
\begin{equation*}
  \min_{\bx\in\RR^n}\;\frac1m\sum_{i=1}^m
  \bigl|\langle\mathbf a_i,\bx\rangle^2-b_i\bigr|.
\end{equation*}
This is~Problem \eqref{prob:main} with $f\equiv0$, $h=\frac1m\|\cdot\|_1$, and
$\calA(\bx)_i=\langle\mathbf a_i,\bx\rangle^2-b_i$; the $\ell_1$ loss makes the
model robust to outliers.

\textbf{Robust low-rank matrix recovery}~\cite{charisopoulos2021low}.
Given a linear measurement map $\mathcal M\colon\RR^{d_1\times d_2}\to\RR^m$ and
observations $\mathbf b$ containing outliers, a rank-$r$ matrix is recovered in
factorized form via
\begin{equation*}
  \min_{\mathbf U\in\RR^{d_1\times r},\,\mathbf V\in\RR^{d_2\times r}}
  \frac1m\bigl\|\mathcal M(\mathbf U\mathbf V^\top)-\mathbf b\bigr\|_1.
\end{equation*}
Here $\bx=(\mathbf U,\mathbf V)$, $h=\frac1m\|\cdot\|_1$, and
$\calA(\bx)=\mathcal M(\mathbf U\mathbf V^\top)-\mathbf b$ is bilinear in
$(\mathbf U,\mathbf V)$, so the model again has the form~Problem \eqref{prob:main}.

\textbf{Minimizing a pointwise maximum}~\cite{zhang2020smoothed}.
For smooth losses $f_1,\dots,f_m$, such as the losses on $m$ data subgroups, one solves
\begin{equation*}
  \min_{\bx\in\RR^n}\;\max_{1\le i\le m} f_i(\bx).
\end{equation*}
This is~Problem \eqref{prob:main} with $h(u)=\max_{1\le i\le m}u_i$ and
$\calA(\bx)=(f_1(\bx),\dots,f_m(\bx))$; a smooth regularizer, if present, enters
through $f$, and the proximal mapping of $h$ reduces to a projection onto the unit
simplex.

\subsection{Related Work}

A standard way to handle the nonsmooth composition in Problem~\eqref{prob:main}
is to introduce the auxiliary variable $\by=\calA(\bx)$ and rewrite it as the
equality-constrained problem
\begin{equation}\label{prob:split}
  \min_{\bx,\by}\; f(\bx)+h(\by)\quad\text{s.t.}\quad \calA(\bx)=\by.
\end{equation}
Because $h\circ\calA$ is neither globally smooth nor directly prox-friendly,
methods that solve inner subproblems---such as prox-linear
schemes~\cite{Drusvyatskiy2019}, quadratic penalty methods
\cite{kong2019complexity,lin2022complexity}, and inexact or proximal augmented
Lagrangian methods (ALMs)~\cite{sahin2019inexact,li2021rate,xie2021proximal}---incur
inner-loop overhead and termination-criteria tuning. For large-scale problems,
single-loop first-order methods, which perform one explicit gradient or
proximal step per variable at each iteration without an inner solver, are
therefore attractive. Yet without additional regularity assumptions, existing
single-loop methods for Problem~\eqref{prob:main} stop at $\mathcal O(\epsilon^{-3})$
complexity and require an algorithm parameter to scale with $1/\epsilon$.
Moreau-envelope smoothing~\cite{bohm2021variable} sets the smoothing parameter
to $\Theta(\epsilon)$ to control the approximation bias (for linear $\calA$),
which inflates the surrogate gradient's Lipschitz constant to $\Theta(\epsilon^{-1})$.
Similarly, single-loop penalty and linearized ALM/ADMM schemes acting
on~\eqref{prob:split} drive the penalty or dual-perturbation parameter to
$\Theta(\epsilon^{-1})$ to control feasibility
\cite{hajinezhad2019perturbed,luSingleLoopGradientDescent2022,hallak2023adaptive,alacaoglu2024complexity,sun2024dual_descent_alm_admm,elbourkhissi2025perturbed,shi2025momentum};
a large penalty likewise inflates the primal smoothness constant and shrinks
the primal stepsize to $\Theta(\epsilon)$. Fixed-penalty $\mathcal O(\epsilon^{-2})$
guarantees are available under structural regularity conditions---such as
polyhedral error bounds for affine constraints~\cite{zhang2022global} or
uniform LICQ for smooth nonlinear constraints
\cite{elbourkhissi2025linearized,liu2026fixed}---but do not apply to the
general nonsmooth composite problem~\eqref{prob:main}, where we impose no such
assumptions.

Eliminating $\by$ from the Lagrangian of~\eqref{prob:split} yields the
equivalent nonconvex--linear minimax formulation
\begin{equation}\label{eq:ncl}
  \min_{\bx\in\RR^n}\max_{\bq\in\RR^m}\;
  \calL(\bx,\bq) = \Phi(\bx,\bq)-h^*(\bq), \qquad
  \Phi(\bx,\bq):=f(\bx)+\langle\bq,\calA(\bx)\rangle,
\end{equation}
whose game-stationary (GS) points match KKT points of~\eqref{prob:split}.
Likewise, for any fixed penalty $\rho>0$, minimizing the augmented Lagrangian
of~\eqref{prob:split} over $\by$ produces a smooth nonconvex--concave (NC--C)
minimax problem on which single-loop linearized ALMs act as gradient
descent--ascent (GDA) schemes (Section~\ref{subsec:reduced-al}). For smooth
NC--C minimax problems, two-timescale GDA requires $\mathcal O(\epsilon^{-6})$
iterations~\cite{lin2020gradient}. Adding primal smoothing (Smoothed
GDA~\cite{zhang2020smoothed}) or dual perturbation (Perturbed GDA /
AGP~\cite{lu2020hybrid,xu2023unified}) improves the complexity to
$\mathcal O(\epsilon^{-4})$, and combining both in Perturbed Smoothed GDA
achieves $\mathcal O(\epsilon^{-3})$ with a dual regularization parameter of
order $\epsilon^{-1}$; Li et al.~\cite{li2026smoothing} establish matching
algorithm-specific lower bounds of $\Omega(\epsilon^{-4})$ and
$\Omega(\epsilon^{-3})$ for Smoothed GDA and Perturbed Smoothed GDA,
respectively. Single-loop methods tailored to nonconvex--linear structure
likewise attain $\mathcal O(\epsilon^{-3})$ via $\epsilon$-dependent dual
regularization~\cite{pan2021efficient,he2024approximation,xu2026rada}.

Across all of these $\mathcal O(\epsilon^{-3})$ single-loop guarantees, some
parameter grows with $1/\epsilon$---either the penalty parameter itself or the
dual regularization parameter. In practice, however, a primary appeal of ALMs
over penalty methods is keeping the penalty parameter fixed so that the primal
smoothness constant does not blow up as $\epsilon\to0$. Because existing lower
bounds for single-loop methods apply only to individual algorithms under fixed
parameter rules~\cite{mahdavinia2022tight,li2026smoothing} (whereas
problem-class bounds for smooth nonconvex or affinely constrained problems do
not address this trade-off~\cite{carmon2020lower,liu2026lower}), it remains
unclear whether growing a parameter is fundamentally unavoidable for
single-loop linearized ALMs, and if so, whether the penalty itself must bear
this growth. This motivates our first question:
\begin{quote}
  \textbf{(Q1)}
  \emph{Without additional regularity assumptions, must a single-loop
  linearized ALM let some parameter grow with $1/\epsilon$ to reach
  $\mathcal O(\epsilon^{-3})$ complexity? If so, must this growth be carried by
  the penalty parameter?}
\end{quote}

In this paper, we provide a precise characterization addressing \textbf{(Q1)}.
We study LRALM, a family $\calC_\theta$ of single-loop linearized ALMs that
perform one GDA step per iteration on the reduced augmented Lagrangian with a
fixed penalty $\rho$, a dual regularization parameter $s$, and primal and dual
centers updated with constant center stepsizes $(\alpha,\beta)\in[0,1]^2$
(Section~\ref{sec:algorithm-framework}). The family $\calC_\theta$ subsumes
Smoothed GDA and Perturbed Smoothed GDA as special cases, and every member
shares the effective penalty $\rho+s$ (Lemma~\ref{lem:effective-penalty}). We
answer the first part of \textbf{(Q1)} affirmatively and the second part
negatively. Whenever $\rho+s\le c_1/\epsilon$, every choice of constant center
stepsizes in $\calC_\theta$ requires $\Omega(\epsilon^{-4})$ iterations to find
an $\epsilon$-KKT point (Theorem~\ref{thm:lower-bound}), matching an
$\mathcal O(\epsilon^{-4})$ upper bound in the family
(Theorem~\ref{thm:upper-beta1}). In contrast, keeping $\rho$ fixed and setting
$s=\Theta(\epsilon^{-1})$ attains $\mathcal O(\epsilon^{-3})$ within
$\calC_\theta$ (Theorem~\ref{thm:upper-beta0}), while every choice of constant
center stepsizes in $\calC_\theta$ with fixed $\rho$ and
$s=\Theta(\epsilon^{-1})$ still requires $\Omega(\epsilon^{-3})$ iterations
(Theorem~\ref{thm:lower-bound-fixed-rho}). Hence some parameter must grow to
order $1/\epsilon$, the dual regularization $s$ alone can carry this growth
while $\rho$ stays fixed, and $\Theta(\epsilon^{-3})$ is tight throughout
$\calC_\theta$.

\begin{table}[tbp]
\centering
\begin{threeparttable}
\caption{Comparison of complexity bounds (iterations for single-loop methods;
first-order oracle calls for multi-loop methods) for finding an $\epsilon$-KKT
point (ALM-based methods and ours) or an $\epsilon$-GS point (smooth NC--C
minimax methods).}
\label{tab:known-results}
\scriptsize
\setlength{\tabcolsep}{2.4pt}
\renewcommand{\arraystretch}{1.12}
\setlength{\aboverulesep}{0pt}
\setlength{\belowrulesep}{0pt}
\begin{tabular}{lccccccc}
\toprule
\multirow{2}{*}[-1.2ex]{Method} & \multirow{2}{*}[-1.2ex]{Single-loop}
 & \multirow{2}{*}[-1.2ex]{Nonlinear} & \multirow{2}{*}[-1.2ex]{CQ-free}
 & \multicolumn{3}{c}{Parameter scale}
 & \multirow{2}{*}[-1.2ex]{Complexity} \\
\cmidrule(lr){5-7}
 & & & & Penalty $\rho$ & Dual reg.\ $s$
 & Primal prox.\ $\gamma$ & \\
\midrule
\multicolumn{8}{l}{\textit{ALM-based methods}} \\
PProx-PDA~\cite{hajinezhad2019perturbed}
 & \checkmark & -- & \checkmark & $\Theta(\epsilon^{-1})$ & $\Theta(\epsilon^{-1})$ & --
 & $\mathcal O(\epsilon^{-4})$ \\
iALM~\cite{li2021rate,kong2023iteration}
 & -- & -- & -- & $\Theta(\epsilon^{-1})$ & -- & $\Theta(1)$
 & $\tilde{\mathcal O}(\epsilon^{-5/2})$ \\
S-prox-ALM~\cite{zhang2022global}
 & \checkmark & -- & \checkmark & $\Theta(1)$ & -- & $\Theta(1)$
 & $\mathcal O(\epsilon^{-2})$ \\
GDPA~\cite{luSingleLoopGradientDescent2022,sun2024dual_descent_alm_admm}
 & \checkmark & \checkmark & -- & $\Theta(\epsilon^{-1})$ & $\Theta(\epsilon^{-1})$ & --
 & $\mathcal O(\epsilon^{-3})$ \\
NR-LALM~\cite{liu2026fixed}
 & \checkmark & \checkmark & -- & $\Theta(1)$ & -- & --
 & $\mathcal O(\epsilon^{-2})$ \\
\midrule
\multicolumn{8}{l}{\textit{Smooth NC--C minimax}\tnote{a}} \\
Smoothed GDA~\cite{zhang2020smoothed,li2026smoothing}
 & \checkmark & \checkmark & \checkmark & -- & -- & $\Theta(1)$
 & $\Theta(\epsilon^{-4})$ \\
Perturbed Smoothed GDA~\cite{li2026smoothing}
 & \checkmark & \checkmark & \checkmark & -- & \cellcolor{gray!25}$\Theta(\epsilon^{-1})$ & $\Theta(1)$
 & \cellcolor{gray!25}$\Theta(\epsilon^{-3})$ \\
Alg.~1 of~\cite{zhang2026matching}\tnote{b}
 & \checkmark & \checkmark & \checkmark & -- & $\Theta(\epsilon^{-1})$ & $\Theta(1)$
 & $\mathcal O(\epsilon^{-5/2})$ \\
Tracked-FOAM~\cite{pan2026optimal}
 & -- & \checkmark & \checkmark & -- & $\Theta(\epsilon^{-1})$ & $\Theta(1)$
 & $\Theta(\epsilon^{-5/2})$ \\
\midrule
\multicolumn{8}{l}{\textit{Ours}} \\
\phantom{\textbf{LRALM} $\boldsymbol{\calC_\theta}$} (Thms.~\ref{thm:upper-beta1}, \ref{thm:lower-bound})
 & \checkmark & \checkmark & \checkmark & \cellcolor{gray!25}$\Theta(1)$ & \cellcolor{gray!25}$\Theta(1)$ & $\Theta(1)$
 & \cellcolor{gray!25}$\Theta(\epsilon^{-4})$ \\
\multirow{-2}{*}{\textbf{LRALM} $\boldsymbol{\calC_\theta}$} (Thms.~\ref{thm:upper-beta0}, \ref{thm:lower-bound-fixed-rho})
 & \checkmark & \checkmark & \checkmark & \cellcolor{gray!25}$\Theta(1)$ & \cellcolor{gray!25}$\Theta(\epsilon^{-1})$ & $\Theta(1)$
 & \cellcolor{gray!25}$\Theta(\epsilon^{-3})$ \\
\textbf{Fast LRALM} (Alg.~\ref{alg:b2-prime})\tnote{c}
 & \checkmark & \checkmark & \checkmark & \cellcolor{gray!25}$\Theta(1)$ & \cellcolor{gray!25}$\Theta(\epsilon^{-1})$ & $\Theta(1)$
 & \cellcolor{gray!25}$\mathcal O(\epsilon^{-5/2})$ \\
\bottomrule
\end{tabular}
\begin{tablenotes}[flushleft]
\scriptsize
\item[] ``Nonlinear'': nonlinear constraints/couplings allowed; ``CQ-free'': no
constraint qualification required. Gray shading highlights the scales of
$\rho,s$ and the resulting complexity.
\item[a] Lower bounds for Smoothed GDA and Perturbed Smoothed GDA are
algorithm-specific; that for Tracked-FOAM holds over the smooth NC--C class.
\item[b] Applied to $\mathcal D_\rho$ on the dual ball of radius $2L_h$, this
method also yields $\mathcal O(\epsilon^{-5/2})$ for Problem~\eqref{prob:main}.
\item[c] Whether $\mathcal O(\epsilon^{-5/2})$ is tight for
Problem~\eqref{prob:main} remains open.
\end{tablenotes}
\end{threeparttable}
\end{table}

Our resolution of \textbf{(Q1)} shows that $\Theta(\epsilon^{-3})$ is an
unimprovable barrier for the LRALM family $\calC_\theta$. Faster rates of order
$\epsilon^{-5/2}$ are nevertheless known in related settings. For smooth NC--C
minimax problems, the optimal deterministic first-order complexity for
$\epsilon$-GS points is $\Theta(\epsilon^{-5/2})$~\cite{pan2026optimal}
(reflecting the $\Omega(\sqrt{\kappa}\,\epsilon^{-2})$ lower bound for
nonconvex--strongly-concave problems with dual condition number
$\kappa=\Theta(\epsilon^{-1})$~\cite{zhang2021complexity,li2021complexity}),
attained by multi-loop methods~\cite{pan2026optimal} and, recently, by a
single-loop projected extragradient scheme with dual gradient
averaging~\cite{zhang2026matching}. For nonsmooth composite minimax and
constrained problems, however, existing $\epsilon^{-5/2}$-order methods all
rely on multi-loop architectures: composite minimax methods attain
$\mathcal O(\epsilon^{-5/2})$ via nested proximal-point or anchored
schemes~\cite{kong2021accelerated,ke2026efficient}, and inexact ALM and penalty
methods for affinely constrained problems attain
$\widetilde{\mathcal O}(\epsilon^{-5/2})$ using accelerated inner
solvers~\cite{li2021rate,lin2022complexity,kong2023iteration}. This contrast
between the $\Theta(\epsilon^{-3})$ barrier of $\calC_\theta$ and the known
$\mathcal O(\epsilon^{-5/2})$ complexity motivates our second question:
\begin{quote}
  \textbf{(Q2)}
  \emph{Can a single-loop linearized ALM with a fixed, $\epsilon$-independent
  penalty parameter break the $\Theta(\epsilon^{-3})$ barrier of $\calC_\theta$
  and attain $\mathcal O(\epsilon^{-5/2})$ complexity for
  Problem~\eqref{prob:main}?}
\end{quote}

To answer \textbf{(Q2)}, we introduce Fast LRALM (Algorithm~\ref{alg:b2-prime}
in Section~\ref{sec:acceleration}). A key structural consequence of our
analysis of \textbf{(Q1)} is that, with $\rho$ fixed and
$s=\Theta(\epsilon^{-1})$, fixing the dual center at zero ($\beta=0$) incurs no
loss in complexity order while making the regularized dual objective fixed,
smooth, and $s^{-1}$-strongly concave on all of $\RR^m$. Stepping outside
$\calC_\theta$, Fast LRALM adds Nesterov extrapolation directly to the dual
multiplier step while preserving a single explicit primal--dual update per
iteration, with no inner solvers, dual projections, or extragradient steps.
We prove that Fast LRALM finds an $\epsilon$-KKT point of
Problem~\eqref{prob:main} in $\mathcal O(\epsilon^{-5/2})$ iterations with a
fixed penalty $\rho$ and $s=\Theta(\epsilon^{-1})$
(Theorem~\ref{thm:b2-prime-complexity}), improving on the
$\Theta(\epsilon^{-3})$ complexity of $\calC_\theta$ and matching the
$\mathcal O(\epsilon^{-5/2})$ rate of multi-loop methods (and
of~\cite{zhang2026matching} applied to the reduced augmented Lagrangian on a
bounded dual ball) via a simpler single-loop ALM update. Because the dual step
is unconstrained and extrapolated, establishing uniform multiplier boundedness
without projection and controlling the $\sqrt{\kappa}$-amplified tracking error
under a moving primal center require a dedicated Lyapunov analysis, developed
in Section~\ref{sec:acceleration}. Table~\ref{tab:known-results} compares our
guarantees and parameter scales with representative methods.

\paragraph{Organization}
Section~\ref{sec:preliminaries} introduces the problem class, the KKT
criterion, the oracle model, and the reduced augmented Lagrangian together with
its regularized potential.
Section~\ref{sec:algorithm-framework} introduces the LRALM family and the
complexity measure, shows that every member of the family has the effective
penalty $\rho+s$, and establishes the lower and upper bounds.
Section~\ref{sec:acceleration} presents Fast LRALM and its iteration
complexity. Section~\ref{sec:experiments} reports numerical experiments, and
Section~\ref{sec:conclusion} discusses open questions. The hard instances
and the remaining proofs are given in the appendices.

\paragraph{Notation}
We use $\langle\cdot,\cdot\rangle$ and $\|\cdot\|$ for the
Euclidean inner product and norm; for a matrix, $\|\cdot\|$ is the spectral
norm. For $\calA=(A_1,\ldots,A_m)^\top$, the Jacobian
$J_{\calA}(\bx)\in\RR^{m\times n}$ has $i$th row $\nabla A_i(\bx)^\top$. For a
nonempty closed convex set $C$, $\dist(\bu,C):=\min_{\bv\in C}\|\bu-\bv\|$, and
$\Pi_C(\bu)$ is the Euclidean projection of $\bu$ onto $C$. When convenient, we
identify $\RR^2$ with $\mathbb C$, so that $\|p\|=|p|$ and
$\langle p,q\rangle=\operatorname{Re}(\bar pq)$.
% Throughout the paper, $\partial h$ denotes the limiting
% subdifferential of $h$.

For a proper closed convex function \(h:\RR^m\to\RR\cup\{+\infty\}\), its
subdifferential, convex conjugate, and proximal mapping are
\[
    \partial h(\by):=\{\bv:\ h(\bw)\ge h(\by)+\langle\bv,\bw-\by\rangle
    \ \text{for all }\bw\},
    \qquad
    h^*(\bq):=\sup_{\by\in\RR^m}\{\langle\bq,\by\rangle-h(\by)\},
\]
\begin{equation}
    \prox_{t h}(\bu)
    :=\argmin_{\by\in\RR^m}
       \Bigl\{h(\by)+\frac{1}{2t}\|\by-\bu\|^2\Bigr\},
    \qquad t>0;
    \label{eq:prox-definition}
\end{equation}
see~\cite{rockafellar1970convex}. We write
\(\dom\partial h^*:=\{\bq:\partial h^*(\bq)\neq\varnothing\}\).
